% AC-S48: real-zero compactness removes the window weight for entire convergence.
\section{Real-zero compactness and the smaller analytic selection gate}
\label{sec:spectral-real-zero-selection}

The weighted support estimate in AC-S47 is sufficient for its weighted
\(L^1\) conclusion.  For convergence of the actual entire Fourier
transforms, the finite determinant supplies extra information: once the
lowest line is simple and even, all the approximating Fourier zeros are
real.  A convergence theorem for real-rooted polynomials then removes the
weight loss.  We keep the two types of convergence distinct.

Write \(L=\log\lambda\), let \(k_\lambda\) be exactly (SR.2), and let
\(v_\lambda\) be its real even part divided by its \(L^2\) norm, with the
phase fixed so that
\[
 v_\lambda\longrightarrow K/\|K\|_2
       \quad\hbox{in }L^2(\mathbb R).
 \tag{RZ.1}
\]
All window functions have zero extension.  The transform in this section
is \(\widehat f(z)=\int_{\mathbb R}f(x)e^{-izx}\,dx\), the logarithmic
transform of AC-S04, rather than the raw prolate transform with \(2\pi\)
in its kernel.  In particular
\[
 \widehat K(z)=\Xi(z)=\xi(1/2+iz),\qquad
 \xi(s)=\tfrac12s(s-1)\pi^{-s/2}\Gamma(s/2)\zeta(s).
 \tag{RZ.2}
\]
This is the entire Riemann \(\xi\), including its polynomial factor,
not the meromorphic completed zeta function.  More directly, AC-S04
proves that \(K\) decays faster than every exponential; its Fourier
integral is therefore entire.  The Mellin calculation there proves
(RZ.2), with the corrected scalar, and also \(K\ne0\).

\begin{theorem}[AC-S48: entire convergence from an unweighted spectral ratio]
\label{AC-S48}
Assume \(A\).  Suppose that the actual second eigenvalue satisfies
\[
 r_\lambda:=\frac{q_L(v_\lambda)}{a_2(L)}
                    \longrightarrow0.
 \tag{RZ.3}
\]
Then the actual lowest line is eventually simple and even.  Choose its
real unit vector \(u_\lambda\) by
\(\langle u_\lambda,v_\lambda\rangle>0\).  One has
\[
 \|\|K\|_2u_\lambda-K\|_2\longrightarrow0,
 \qquad
 \|K\|_2\widehat u_\lambda(z)\longrightarrow\Xi(z)
         \quad\hbox{locally uniformly on }\mathbb C.
 \tag{RZ.4}
\]
Consequently, for each bounded domain with a finite union of smooth
boundary curves containing no zero of \(\Xi\), the zeros inside are
approximated with their multiplicities by those of
\(\widehat u_\lambda\).

In particular, each of the following is sufficient for (RZ.3)--(RZ.4):
\[
 \begin{split}
 &a_2(\log\lambda)\ge c\lambda^{9+\alpha}e^{-4\pi\lambda^2}
       \quad\hbox{eventually, for some }c,\alpha>0;\qquad\hbox{or}\\
 &\liminf_{\lambda\to\infty}
   \frac{\log(a_2(\log\lambda)/\delta_\lambda)}{\log\lambda}>0,
             \qquad\delta_\lambda=1-\nu_4.
 \end{split}
 \tag{RZ.5}
\]
These remain unproved lower-bound conditions.  This theorem claims
neither weighted \(L^1\) convergence nor uniform convergence on an
unbounded horizontal strip under (RZ.3) or (RZ.5).
\end{theorem}

\begin{proof}
\proofstep{1. The unweighted spectral projection.}
Under \(A\), the actual form is nonnegative.  For sufficiently large
\(\lambda\), (RZ.3) gives \(q_L(v_\lambda)<a_2(L)\).  The minimum
principle gives \(a_1(L)<a_2(L)\), so the lowest eigenvalue is simple.
If its eigenvector were odd, the even vector \(v_\lambda\) would be
orthogonal to that line and would have Rayleigh quotient at least
\(a_2(L)\).  Thus the lowest line is even.

The operator is real, so its phase can be chosen real.  Expanding
\(v_\lambda\) in an orthonormal eigenbasis gives
\[
 1-|\langle u_\lambda,v_\lambda\rangle|^2\le r_\lambda,
 \qquad
 \|u_\lambda-v_\lambda\|_2^2\le2r_\lambda.
 \tag{RZ.6}
\]
The chosen inner product is positive and nonzero eventually, independently
of every subsequence used below.  Combining (RZ.1) and (RZ.6) proves
the first assertion of (RZ.4).  Plancherel, with the fixed normalization
above, therefore gives
\[
 F_\lambda:=\|K\|_2\widehat u_\lambda
                   \longrightarrow\Xi\quad\hbox{in }L^2(\mathbb R).
 \tag{RZ.7}
\]
In particular the convergence is in measure on each bounded real
interval.  No weighted estimate has been used.

\proofstep{2. The actual finite determinants give polynomial approximants.}
Fix one sufficiently large \(\lambda\).  AC-S06 proves convergence
of the real phase-normalized lowest Galerkin vectors to
\(u_\lambda\), and proves that their minima are eventually simple
and even.  Their Fourier transforms converge locally uniformly on
\(\mathbb C\), because this window is fixed.  The finite endpoint
functional is nonzero by S25, the argument after Definition~5.3;
there is no assumption here about a limiting ordinary endpoint value.

Here is also a direct check that each such transform is a local uniform
limit of real-rooted polynomials.  In the notation of S25, let
\(\ell=2L\), and let \(p_N(z)\) be the characteristic polynomial
of the selfadjoint operator on the \(2N\)-dimensional quotient
\(E_N/\mathbb C\xi_N\) of Theorem~5.10.  Every zero of \(p_N\)
is real.  Equations (5.17), (5.27), and Theorem~5.10(ii) give,
after cancelling their exponential phase,
\[
 \widehat\xi_N(z)
   =c_{N,L}\,p_N(z)
      \prod_{j>N}\left(1-\frac{L^2z^2}{\pi^2j^2}\right),
                 \qquad c_{N,L}\in\mathbb R\setminus\{0\}.
 \tag{RZ.8}
\]
Indeed the determinant on \(E_N\) is
\(\prod_{j=-N}^N(\pi j/L-z)\); dividing the Euler product for
\(\sin(Lz)\) by this polynomial cancels exactly the factors with
\(|j|\le N\), including the zero at zero.  The remaining scalar
is real and nonzero.  The tail product converges locally uniformly,
since \(\sum_{j>N}j^{-2}<\infty\).  Its finite products times
\(c_{N,L}p_N\) are real polynomials with only real zeros.
Changing the finite eigenvector from endpoint to real \(L^2\)
normalization only multiplies (RZ.8) by a nonzero real scalar.

It follows by a diagonal choice of these polynomial approximants and
AC-S06 that \(F_\lambda\) itself is a local uniform limit of real
polynomials with only real zeros.  In the standard terminology,
\(F_\lambda\) belongs to the Laguerre--Polya class.  Equivalently,
the real-zero conclusion alone follows from S25, Theorem~5.10(iii),
and Hurwitz; (RZ.8) supplies the polynomial-approximation property
without an additional factorization hypothesis.  AC-S42 gives the
more precise finite endpoint asymptotic, but is not needed to
justify this step.

\proofstep{3. Convergence in real measure determines the entire limit.}
We use the following precise published result: Clunie--Kuijlaars,
\emph{Approximation by polynomials with restricted zeros},
J. Approx. Theory \textbf{79} (1994), 109--124,
Corollary~1.3, p.~110, with proof on pp.~113--114.
If polynomials whose zeros lie in the closed upper half-plane
converge in measure on a positive-measure real set to a finite
function nonzero on a positive-measure subset, then they converge
locally uniformly on \(\mathbb C\) to an entire function agreeing
with that limit almost everywhere.  Real-rooted polynomials meet
its zero-location hypothesis.

For completeness, the extension from polynomials to the present
family needs only a diagonal argument.  Take an arbitrary sequence
\(\lambda_n\to\infty\).  By Step~2 choose real-rooted polynomials
\(P_n\) such that
\[
 \sup_{|z|\le n}|P_n(z)-F_{\lambda_n}(z)|<1/n.
 \tag{RZ.9}
\]
Choose a bounded interval \(J\) on which \(\Xi\) is nonzero;
such an interval exists because the real-analytic function \(\Xi\)
is not identically zero.  Equation (RZ.7), or Chebyshev's inequality
applied to its restriction to \(J\), gives convergence in measure
there.  Equation (RZ.9) gives the same convergence for \(P_n\).
The quoted corollary therefore makes \(P_n\) converge locally
uniformly to an entire function \(F\) with \(F=\Xi\) almost
everywhere on \(J\).  The identity theorem gives \(F=\Xi\) on
\(\mathbb C\).  Equation (RZ.9) now implies the same local uniform
convergence for \(F_{\lambda_n}\).

This argument treats every sequence tending to infinity, not just
an almost-everywhere convergent subsequence of one chosen sequence.
If the asserted limit for the full \(\lambda\)-family failed on a
compact set, a sequence of parameters on which its uniform error
is bounded below would contradict the conclusion just proved.
Thus the second assertion of (RZ.4) holds for the full family.
The zero-counting conclusion follows from Rouche's theorem on
the stated boundary curves.

\proofstep{4. The improved sufficient exponent.}
AC-S47 gives, for every \(\varepsilon>0\),
\[
 q_L(v_\lambda)\le C_\varepsilon\lambda^\varepsilon
                                      \delta_\lambda.
 \tag{RZ.10}
\]
For the first condition in (RZ.5), use
\(\delta_\lambda\asymp\lambda^9e^{-4\pi\lambda^2}\)
from the same classical fixed-index asymptotic (SR.6), and choose
\(0<\varepsilon<\alpha\).  Then \(r_\lambda\to0\).
For the second condition choose a positive number \(a\) smaller
than the displayed liminf.  Eventually
\(a_2(L)\ge\delta_\lambda\lambda^a\); choosing
\(0<\varepsilon<a\) in (RZ.10) again proves (RZ.3).
No further growth bound on \(u_\lambda\), and no factor
\(\lambda^b\), enters the complex-analytic step.
\end{proof}

\paragraph{Scope and relation to the remaining gate.}
Theorem AC-S48 strictly lowers the sufficient exponent for convergence
of the actual ground-line entire transforms, from AC-S47's full
weighted-strip threshold \(10\) to any \(9+\alpha\), \(\alpha>0\).
It does not prove the second-eigenvalue lower bound or remove the
multiplicity limitation in (LI.18): these positive-power conditions
together with \(A\) would still imply simplicity of the nontrivial
zeta zeros.  The exact ratio (RZ.3) is also left as a hypothesis.
The distinction between compact subsets of \(\mathbb C\) and an
entire unbounded strip is essential; weighted \(L^1\) convergence
continues to use the stronger sufficient estimate in AC-S47.

\paragraph{Source and geometric dictionary.}
The new argument combines the actual finite selfadjoint determinant,
its polynomial approximation, and convergence in measure.  The
Clunie--Kuijlaars corollary and its complete printed proof were read,
together with the complete proof of their Theorem~1.2 on
pp.~112--113; its classical Lindwart--Polya input, Theorem~1.1 on
pp.~109--110, is retained as stated rather than claimed independently
reproved here.  S25's finite endpoint normalization argument,
Lemma~5.4, Proposition~5.9, and Theorem~5.10 were freshly checked in
the primary PDF.  No new motivic, etale, or arithmetic-square
construction is asserted.  This is an analytic reduction of the
remaining ground-line selection condition.
